% Held-out identities and CCTV labels: which classifier to trust, and when to
% say "not sure". Every number is an \attr{key} macro that make_figures.py
% writes to figures/attr_numbers.tex from results/attributes.json,
% label_check_before.json, identity_gap.json and rq1_gap.json. A missing
% number prints as ?? in bold.

\providecommand{\attr}[1]{%
  \ifcsname attr/#1\endcsname\csname attr/#1\endcsname\else\textbf{??}\fi}
\InputIfFileExists{figures/attr_numbers.tex}{}{}

\subsection{Held-out Identities and CCTV Labels}
\label{sec:cctv}
The PETA figures above come from image-level splits, in which the same person can
appear in training and test. We therefore re-scored both released models on PETA
images they never trained on, and on a 9.7-minute recording from a ceiling camera
in a lecture room with about forty seated students, viewed from above and behind.
That recording was labelled by an AI annotator: gender and age group for each of
its tracks, from contact sheets of the crops the pipeline classified, with people
it could not judge (two people in the crops, only hands visible) marked unsure.
These labels are not human ground truth; they compare models on the same people,
and are to be confirmed by two human annotators.

\emph{What went wrong.} Before any change, the pipeline's per-person gender label
was right \attr{before/gender} of the time over \attr{before/gender/n} labelled
people, and the age group \attr{before/age}. Looking at the errors showed four
causes. All of a person's votes came from their first few seconds, so one awkward
angle decided the label. Track numbers moved between seated neighbours: the same
person's box overlaps its previous box by an IoU of at least 0.5 in
\attr{gap/same05}\% of samples, but \attr{gap/neighbour03}\% of overlapping
neighbour pairs also clear the old threshold of 0.3, while only
\attr{gap/neighbour05}\% reach 0.5. Some crops showed two people. And both
classifiers were weak on seated, top-down views.

\emph{Which classifier.} Table~\ref{tab:cctv} compares the fine-tuned ViT, CLIP
ViT-B/32~\cite{radford2021clip} used zero-shot (asked whether a crop is a CCTV image
of a man or of a woman, with no training), and the average of the two. The ViT is best on PETA's
small crops and CLIP on the 1080p lecture crops; their average is the best on
both, with a balanced accuracy of \attr{gender/cuhk/ensemble/bal} on CUHK against
\attr{gender/cuhk/vit/bal} for the ViT alone, and \attr{gender/room/ensemble/acc}
accuracy on the lecture room against \attr{gender/room/vit/acc}. For age, CLIP
zero-shot with coarse prompts (a young person, an adult, a middle-aged person, an
old person) reaches \attr{age/peta/clip/bal} balanced accuracy over PETA's four
groups, against \attr{age/cuhk/svm/bal} for the SVMs on CUHK. Logistic heads
fitted on CLIP features of PETA's training split did worse on the lecture room
than zero-shot: PETA's small, standing pedestrians do not look like CCTV of
seated people. Almost every labelled age in the lecture room is 16 to 30, so the
age results there only test that group.

\begin{table}[t]
\caption{Gender and age per person. PETA rows use human labels on images the
models never trained on; lecture-room rows use the AI annotator's labels.}
\label{tab:cctv}
\centering
\begin{tabular}{@{}llcc@{}}
\toprule
Task & Model & PETA & Lecture room \\
\midrule
Gender & Fine-tuned ViT      & \attr{gender/cuhk/vit/bal}      & \attr{gender/room/vit/acc} \\
       & CLIP zero-shot      & \attr{gender/cuhk/clip/bal}     & \attr{gender/room/clip/acc} \\
       & Average of the two  & \attr{gender/cuhk/ensemble/bal} & \attr{gender/room/ensemble/acc} \\
Age    & Colour/texture SVMs & \attr{age/cuhk/svm/bal}         & \attr{age/room/svm/acc} \\
       & CLIP zero-shot      & \attr{age/peta/clip/bal}        & \attr{age/room/clip/acc} \\
\midrule
\multicolumn{4}{@{}l}{PETA: balanced accuracy, CUHK test images for gender and the}\\
\multicolumn{4}{@{}l}{ViT; all test identities for CLIP age. Lecture room: accuracy.}\\
\bottomrule
\end{tabular}
\end{table}

\emph{Saying ``not sure''.} Leaving people whose pooled confidence is below a
threshold without a label trades coverage for accuracy. On the lecture room, the
averaged gender model at a threshold of 0.6 labels \attr{gender/room/ensemble/kept0.6}\%
of people with \attr{gender/room/ensemble/acc0.6}\% accuracy, and CLIP age at 0.4
labels \attr{age/room/clip/kept0.4}\% with \attr{age/room/clip/acc0.4}\%. These
are the pipeline's defaults, together with the other fixes of
Section~\ref{sec:pooling}: votes spread over the whole stay, no votes on crops
showing two people, matching at an IoU of 0.5 in crowded rooms, and an unsure
label for a person whose votes disagree between the two halves of their stay.
